\documentclass[11pt,a4paper]{article}
\usepackage[T1]{fontenc}
\usepackage[utf8]{inputenc}
\usepackage{lmodern,microtype}
\usepackage[margin=.85in]{geometry}
\usepackage{amsmath,amssymb,amsthm,mathtools,aliascnt}
\usepackage{booktabs,tabularx,array,enumitem,needspace,float}
\usepackage[section]{placeins}
\usepackage{xcolor,tikz,mdframed}
\usetikzlibrary{arrows.meta,positioning,calc,fit}
\usepackage[hidelinks]{hyperref}
\usepackage[nameinlink,noabbrev]{cleveref}
\usepackage{fancyhdr}
\hypersetup{pdftitle={Sync HotStuff RR v3},pdfsubject={L1-scheduled leaders and generation-wide recovery}}
\setlength{\headheight}{14pt}
\pagestyle{fancy}\fancyhf{}
\lhead{\small Sync HotStuff RR v3}\rhead{\small L1-scheduled leaders}
\cfoot{\thepage}
\setlength{\parindent}{0pt}\setlength{\parskip}{4pt}
\setlength{\emergencystretch}{2em}
\setlist{nosep,leftmargin=*}
\newcolumntype{Y}{>{\raggedright\arraybackslash}X}
\definecolor{rrRed}{HTML}{A51C30}
\definecolor{rrBlue}{HTML}{2563EB}
\definecolor{rrTeal}{HTML}{0F766E}
\definecolor{rrAmber}{HTML}{B45309}
\newcommand{\RR}[1]{{\color{rrRed}#1}}
\newcommand{\QC}{\mathcal C}
\newcommand{\CB}{\mathsf{CB}}
\newcommand{\Hash}{\mathsf H}
\newcommand{\rank}{\operatorname{rank}}
\newcommand{\tagg}[1]{\mathsf{#1}}
\newtheorem{definition}{Definition}[section]
\newaliascnt{assumption}{definition}
\newtheorem{assumption}[assumption]{Assumption}
\aliascntresetthe{assumption}
\newaliascnt{lemma}{definition}
\newtheorem{lemma}[lemma]{Lemma}
\aliascntresetthe{lemma}
\newaliascnt{theorem}{definition}
\newtheorem{theorem}[theorem]{Theorem}
\aliascntresetthe{theorem}
\newaliascnt{proposition}{definition}
\newtheorem{proposition}[proposition]{Proposition}
\aliascntresetthe{proposition}
\crefname{assumption}{assumption}{assumptions}
\Crefname{assumption}{Assumption}{Assumptions}
\crefname{lemma}{lemma}{lemmas}
\Crefname{lemma}{Lemma}{Lemmas}
\crefname{theorem}{theorem}{theorems}
\Crefname{theorem}{Theorem}{Theorems}
\newcounter{protocol}
\newenvironment{protocol}[2]{\refstepcounter{protocol}\label{#1}
\begin{mdframed}[nobreak=true,linewidth=.6pt,skipabove=8pt,skipbelow=8pt,
innertopmargin=7pt,innerbottommargin=7pt]\small
\textbf{Protocol \theprotocol. #2}\par\smallskip}{\end{mdframed}}
\crefname{protocol}{Protocol}{Protocols}
\Crefname{protocol}{Protocol}{Protocols}
\title{Sync HotStuff RR v3\\\large Immediate QC Rotation with Accountable Leader Eligibility}
\author{}\date{11 October 2026}
\begin{document}
\maketitle
\begin{abstract}
We specify an immediate-QC rotation variant of standard-synchronous
Sync HotStuff. A finalized L1 authority determines the leader $L(g,t)$
of each generation and term. A terminal DATA certificate admits the next
term immediately, while the old commit timers continue. A complete quit
proof freezes the whole generation, and every replica recovers to
$(g+1,0)$ using SH's entry locks, NEWVIEW, and anchor re-certification.
Correct replicas that learn certified data through catch-up start the
same $2\Delta$ commit timer without casting a redundant vote. Safety
depends on synchrony and authenticated histories. Eventual liveness
additionally uses an accountable PoS service that permanently excludes
revealed equivocators from future proposer schedules after finite L1
processing and epoch activation. We distinguish this eventual guarantee
from an honest-window progress bound and from optional pacing assumptions.
Normal handoff requires neither predecessor commitment nor a DONE certificate.
\end{abstract}

\section{State machine replication and scope}
\label{sec:smr}
A deterministic state machine executes commands ordered in hash-linked
blocks. Committing a block commits its uncommitted ancestors in order.
\begin{definition}[SMR]
Safety means that correct replicas' committed logs are prefix-comparable
and completed commitments are never withdrawn. Validity means that
committed blocks satisfy the application predicate. Request liveness
means that every well-formed correct-client request remaining admissible
for log inclusion until ordered is eventually committed by every correct
replica.
\end{definition}
Committed-height progress is the intermediate property that every correct
replica eventually commits arbitrarily large heights. Request liveness
also needs a fair input service.

Each ordinary term has a quota $K\ge1$ of fresh batches. A normally
completed term certifies all $K$; a faulty term can produce none and be
replaced. Thus a positive quota does not promise Byzantine cooperation.
The two core transitions are
\[
 (g,t)\xrightarrow{\text{terminal DATA QC}}(g,t+1),
 \qquad
 g\xrightarrow{\text{complete quit proof}}(g+1,0).
\]
The first retains outstanding timers; the second cancels every unfinished
timer in generation $g$. A QC remains distinct from a commitment.
This manuscript extends the standard-synchronous protocol in
Section II of Sync HotStuff~\cite{sync}.

The PoS integration analyzed here changes \emph{proposer eligibility};
the voting committee and its threshold remain fixed. Removing a proposer
from future leader lists is sufficient for the stated exclusion argument.
Changing the voting committee or stake-weighted quorum requires a separate
cross-configuration safety construction and is not implicit in this model.

\section{Model, L1 authority, and PoS exclusion}
\label{sec:model}
\begin{assumption}[Authenticated execution]
\label{ass:network}
The committee has $n=2f+1$ members, at most $f$ of which are permanently
Byzantine. Correct replicas retain state and participate throughout.
Complete correct-to-correct messages, including dependencies and validation,
are processed within $\Delta>0$. Local clocks have equal rates but may
have different offsets. Signatures are unforgeable, $\Hash$ is collision
resistant, and bounded time intervals contain finitely many events.
Received evidence is processed before timer callbacks at an equal timestamp.
Correct replicas share genesis $B_0$, the initial finalized authority,
and enter $(0,0)$ within $\Delta$ of one another.
\end{assumption}

The proposed physical-delay profile may additionally supply
$0<\underline\delta\le\Delta$ as a lower bound on delivered inter-replica
messages; self-processing is local. A correct-to-correct delivery then
lies in $[\underline\delta,\Delta]$. Neither safety nor the eventual PoS
liveness proof relies on a positive lower bound. In particular, a wire-delay
bound alone is not assumed to imply a fresh quorum round trip for every
block. We make the distinction precise in \cref{sec:pacing}.

\begin{assumption}[Unique finalized authority and readiness]
\label{ass:l1}
For every generation $g$, L1 eventually provides one immutable finalized
record $\tagg{AUTH}_g$ identifying a PoS epoch $e(g)$ and a unique schedule
$L(g,t)$ over the fixed committee. Correct replicas verify the same
record, retain it, and immediately relay its complete finality proof,
including while frozen. If $d_r(g)$ is the first time correct replica $r$
has that record available, these readiness times differ by at most $\Delta$.
Records may be published in advance. L1 processing and epoch activation
progress independently of L2 commitment.
\end{assumption}
PoS epoch $e(g)$ and recovery generation $g$ are different counters.
An exclusion changes later authority records; it never rewrites the
leader of a previously admitted slot. If the next record is already
available, recovery has no L1 wait. Otherwise it safely waits for that
record; no fixed local-only recovery bound is claimed for an unbounded
L1 availability delay.

\begin{assumption}[Eventual permanent equivocation exclusion]
\label{ass:exclusion}
Every admissible leader-equivocation proof revealed to a correct replica
is relayed and eventually submitted to the independently live L1 service.
After finite finalization and epoch activation, the exposed identity is
permanently ineligible for future proposer schedules. Precisely, for each
such identity $p$ there is a finite generation index $G_p$ such that
\[
 L(g,t)\ne p\qquad\text{for every }g\ge G_p\text{ and }t\ge0.
\]
Excluded identities cannot reenter under replacement identities, and the
finite fixed Byzantine set is not replenished by new corruptions or joins.
The exclusion predicate covers the signed DATA and NEWVIEW equivocations
defined in \cref{sec:evidence}. BLAME alone is not a slashing proof.
\end{assumption}
This is an enforced control-service property, not an assumption that a
Byzantine party values its deposit. Next-epoch removal must have the
stated prospective effect on newly activated schedules; infinitely reusing
an old snapshot would not satisfy the assumption. The node can remain a
faulty voter under the fixed $f$ bound after losing proposer eligibility.

\begin{assumption}[Schedule fairness and input service]
\label{ass:fairness}
For every $g_0$, some finite $g\ge g_0$ has correct $L(g,0)$. For every
fixed $g$ and $t_0$, some finite $t\ge t_0$ has correct $L(g,t)$.
Correct leaders can produce valid empty blocks.

Correct-client requests are eventually disseminated to all correct
replicas and have a common admission order with finitely many predecessors
per request. Each request fits in a batch. Correct leaders fill fresh
blocks with the oldest pending admissible requests that fit, removing a
request only after commitment. A request already in the proposed ancestry
need not be duplicated. Correct replicas retain and serve data and evidence.
\end{assumption}
Cyclic schedules over eligible identities can supply fairness. The
theorems state this property explicitly rather than deriving it from
unspecified L1 randomness. Safety uses neither fairness nor eventual
exclusion; those are additional liveness assumptions.

\section{Views, certification, and admission}
\label{sec:notation}
Views $v=(g,t)$ are ordered lexicographically. Signatures bind the
protocol and voting-configuration domain, $\tagg{AUTH}_g$, message type,
view, height, block hash, opening identity, and parent where applicable.
A correct replica signs at most one DATA block at each $(g,t,k)$.
A fresh block is $B_k=(x_k,\Hash(B_{k-1}))$.

A DATA QC $\QC_{g,t}(B_k)$ consists of $f+1$ matching signatures from
distinct members. ANCHOR votes are separately tagged; their QC
re-certifies an unchanged inherited block. Both QC types have rank
$(g,t,k)$. A fixed hash order resolves equal-rank distinct blocks when a
highest QC must be selected. Different representations of a certificate
do not create different block values.

For $g>0$, any complete valid quit proof for $g-1$ is an entry witness
$E_g$ for the same target $(g,0)$. The leader opens with
\[
 R_g=\langle\tagg{NEWVIEW},(g,0),A_g,Q_g,E_g\rangle_{L(g,0)},
\]
where $Q_g$ certifies $A_g$ in a strictly older generation. The synthetic
genesis certificate has rank $(-1,0,0)$. Generation zero has the fixed
genesis opening and needs no ANCHOR exchange.

The opening identity is $(g,h(A_g),\Hash(A_g))$. Term $t$ owns
\[
 h(A_g)+tK+1,\ldots,h(A_g)+(t+1)K.
\]
The opening is certified, not assumed committed. Its ANCHOR vote consumes
no quota and starts no timer. A DATA history is admitted by the signed
opening and its ANCHOR QC, except at genesis.

\subsection{Immediate terminal-QC entry}
\label{sec:normal}
The terminal block is $T_{g,t}=B_{h(A_g)+(t+1)K}$ and its DATA QC is
$N_{g,t}$. Receiving $N_{g,t}$ admits $(g,t+1)$ immediately; entering a
later term requires every intervening terminal QC. The new leader can
propose and correct replicas can DATA-vote before the predecessor commits.
Old timers continue. No DONE or post-commit certificate gates this path.

Every admitted proposal and QC carries its complete signed history:
opening, authority and entry witnesses, intervening terminal QCs, and
parent certificates. A cached dependency may be referenced only when it
is available within the complete-message bound. Physical capacity and
the size of these bundles remain part of the execution model and byte cost.

\subsection{Quit evidence}
\label{sec:evidence}
Admissibility is checked against authenticated authority and entry history,
not against a recipient's preferred uncommitted branch. Term zero is
admissible on its entry witness even if its leader remains silent. A
later term needs the preceding terminal-QC chain.
A complete $\CB(g,t)$ is:
\begin{enumerate}[label=(\alph*)]
\item $f+1$ distinct signatures on $\langle\tagg{BLAME},g,t\rangle$
for that admissible term;
\item two distinct valid DATA proposals at the same $(g,t,k)$ signed
by $L(g,t)$, with the histories establishing their admission;
\item two valid signed NEWVIEW messages choosing different opening
blocks for $(g,0)$.
\end{enumerate}
Case (c) includes comparable but different openings: they give different
quota origins. Different encodings or different QCs for the same opening
block are not equivocation. With one opening, the first divergent fresh
height in two incompatible histories belongs to one term and one leader,
giving case (b). Single BLAME shares do not freeze a generation.

A replica freezes a generation at most once. Old-generation proofs are
archival; they cannot freeze the current generation. Future-generation
proofs are buffered until authenticated entry. A complete current-generation
proof for a locally unreached term is actionable when its whole
terminal-QC admission is present. Signed equivocation evidence is also
forwarded to the L1 exclusion service, independently of the local recovery.

\subsection{Vote and catch-up timers}
\label{sec:timers}
There is at most one timer for a DATA certificate key
$(g,t,k,\Hash(B_k))$. Its first start event is either a local DATA vote
or receipt of a valid current-generation DATA QC when the replica has
not already committed that block and has no timer for that key. In the
second case it first compares the signed opening and history with retained
evidence, processing any resulting quit proof before a timer can start.
If still unfrozen, it relays the QC with its signed history and starts a
$2\Delta$ \emph{catch-up timer}, without a redundant vote.

Apply the QC rule to DATA QCs embedded in admitted ancestry, not just the
tip QC. It applies even after passing the block's ordinary term. It
never starts from an ANCHOR QC, in a frozen generation, or for an old
generation. Duplicate messages do not restart or shorten timers.
Both timer types commit only if their generation remains unfrozen.
This rule supplies local finality to a replica that catches up through
QCs without having voted for the older blocks.

\section{Protocol specification}
\label{sec:protocol}
Black text retains SH's basic propose/vote/commit and lock-based recovery
operations. Red text marks composite views, L1 authority, immediate regular
rotation, catch-up timers, generation-wide exit, and the PoS evidence
interface. L1 exclusion does not replace signature or history validation.

\begin{protocol}{alg:steady}{Steady state with immediate QC rotation}
Let $v=(g,t)$ be the active view and $L=L(g,t)$.
\RR{Perform DATA operations only while $g$ is unfrozen.}
\begin{enumerate}[label=\textbf{\arabic*)},itemsep=5pt]
\item \textbf{Propose.} If $r=L$, on obtaining $\QC_v(B_{k-1})$,
broadcast $\langle\tagg{PROPOSE},B_k,v,\QC_v(B_{k-1})\rangle_L$
where $B_k$ extends $B_{k-1}$.
\RR{Bind the opening and authority, attach the complete admission history,
and use only the $K$ fresh slots. The first child after ordinary rotation
uses the preceding term's terminal QC. A recovered initial term uses its
ANCHOR QC; generation zero uses genesis. Propose as soon as justified,
without waiting for the predecessor's commitment. A correct proposer
immediately processes and DATA-votes its own admitted fresh proposal.}
\item \textbf{Vote and catch up.} On a valid proposal with no known leader
equivocation, forward it, multicast a DATA vote, and start a $2\Delta$
commit timer on the first vote.
\RR{On every admitted proposal or QC, first compare current-generation openings and
histories, including earlier terms. A complete fork freezes the generation
before another vote. Catch up through terminal QCs before applying the
vote rule; if already past a block, relay its QC instead of re-voting.
For each newly learned current-generation DATA QC in the admitted history,
start a $2\Delta$ catch-up timer if its key has no timer or commitment.
An ANCHOR QC starts no timer. A certified child is processable even if the
recipient declined its opening's original ANCHOR vote against a higher
uncommitted entry lock.}
\item \textbf{(Non-blocking) Commit and rotation.} On timer expiry, commit
the block and its uncommitted ancestors.
\RR{Process arriving evidence first; a frozen generation cannot commit
through its canceled timers. On terminal $N_{g,t}$, immediately enter
$(g,t+1)$, clear the pending complaint, disable the completed term's
progress listener, and let the new leader propose and replicas vote.
Retain older current-generation timers and equivocation listeners.
No terminal commit evidence or DONE vote is required. Never withdraw
completed commitments.}
\end{enumerate}
\end{protocol}

\begin{protocol}{alg:recovery}{Generation-wide quit and SH recovery}
\begin{enumerate}[label=\textbf{\roman*)},itemsep=5pt]
\item \textbf{Blame and quit.} If fewer than $p$ fresh proposals trigger
votes after $(2p+4)\Delta$ from local view entry, broadcast BLAME.
Gather and broadcast $f+1$ matching signatures, or broadcast signed
leader-equivocation evidence.
\RR{Schedule a separate check for every $p=1,\ldots,K$, only while that
view is active and unfrozen. Count distinct admitted DATA positions
supported by an own vote or valid DATA QC; ANCHOR votes do not count.
Keep at most one pending complaint. Any complete current-generation
proof freezes all of $g$: stop new DATA/ANCHOR votes and normal entries,
cancel all vote-start and catch-up timers, but retain commitments and
continue processing old votes/QCs. Relay the proof immediately and do
not restart this recovery. Send revealed equivocation to L1 for prospective
exclusion; timeout BLAME alone does not identify a slashable double signer.}
\item \textbf{Status and entry.} Wait $\Delta$ after local quit, select a
highest valid QC, fix the entry lock, and report it to the new leader.
\RR{The target is always $(g+1,0)$. At
$e_r=\max\{a_r+\Delta,d_r(g+1)\}$, install the common finalized
$\tagg{AUTH}_{g+1}$, set the lock to the maximum of the retained lock
and known QCs from generations through $g$, and send this QC with full
history and $E_{g+1}$ to $L(g+1,0)$. Enter the new view and reset its
proposal clock. This entry lock stays fixed for its ANCHOR vote.
Late certificates may update retained highest evidence for later recovery.
Process buffered current-generation quit evidence immediately on entry.}
\item \textbf{New-view.} The new leader waits $2\Delta$ from its own
entry and broadcasts a signed NEWVIEW with a highest valid known QC.
\RR{Its justification is from a strictly older generation, or genesis.
Pin one opening block; two distinct openings in this generation are
equivocation. The L1 authority chooses the leader, not the anchor.
No local BA over CBs or STATUS vectors is run.}
\item \textbf{First vote.} If the NEWVIEW QC ranks at least as high as the
entry lock, forward NEWVIEW and vote for its unchanged anchor.
\RR{Use ANCHOR-VOTE in $(g+1,0)$ at most once and start no timer.
On the ANCHOR QC, propose the first fresh child and resume
\cref{alg:steady}. Inspect signed histories for equivocation before any
vote. Old-generation evidence cannot restart this generation's recovery.}
\end{enumerate}
\end{protocol}

\paragraph{Processing discipline.}
The entry-lock comparison is specific to ANCHOR first voting; it is not
a permanent refusal to process a child backed by a valid new-generation
ANCHOR QC. Compare signed openings before assigning their quota origins.
Deduplicate NEWVIEW by $(g,h(A),\Hash(A))$, and DATA proposals/QCs by
type, view, height, and block hash, independently of QC subsets or encodings.
Retain two distinct witnesses at a conflicting slot or opening and relay
the first complete quit proof. Forward each admitted DATA QC as needed
for catch-up even if the same block had a differently ranked earlier QC.
All stated communication bounds concern this admitted protocol traffic,
not arbitrary Byzantine flooding.

\section{Execution diagrams}
\label{sec:figures}
The diagrams distinguish immediate certification-driven handoff, local
timer commitment, L1 authority readiness, and future PoS exclusion.
The network arrows show representative messages; they are not assumptions
about identical arrival times.
\begin{figure}[htbp]\centering
% Native TikZ fragment. Representative messages; time is not drawn to scale.
\begin{tikzpicture}[
  x=1cm,y=-1cm,font=\fontsize{9.5}{11}\selectfont,
  lane/.style={draw=black!25,densely dashed},
  msg/.style={-{Stealth[length=1.8mm]},semithick,draw=rrBlue},
  lab/.style={fill=white,inner sep=2pt,align=center},
  state/.style={draw=rrBlue!65,fill=white,rounded corners=2pt,
    align=center,inner sep=4pt},
  progress/.style={draw=rrTeal,fill=rrTeal!4,rounded corners=2pt,
    align=center,inner sep=5pt},
  phase/.style={anchor=west,font=\fontsize{9.5}{11}\selectfont\bfseries}
]
  \path[use as bounding box] (0,-.5) rectangle (15.8,11.7);
  \node[align=center] at (2.1,0) {Current leader\\$L(g,t)$};
  \node[align=center] at (7.8,0) {Correct replica\\$P$};
  \node[align=center] at (13.5,0) {Successor leader\\$L(g,t+1)$};
  \foreach \x in {2.1,7.8,13.5} \draw[lane] (\x,.5)--(\x,10.8);

  \node[phase] at (.1,.9) {Terminal certification admits the next term immediately};
  \draw[msg] (2.1,1.5)--node[lab] {terminal DATA proposal $T$} (7.8,2.0);
  \draw[msg] (2.1,1.5)--(13.5,2.0);
  \node[state,minimum width=11.8cm] at (7.9,2.7)
    {Vote for $T$ in $(g,t)$; start its $2\Delta$ timer at the local DATA vote.\\
     Broadcast DATA votes; terminal certificate $N_{g,t}=\QC_{g,t}(T)$ forms.};
  \draw[msg] (7.8,3.45)--node[lab] {$N_{g,t}$ + complete history} (13.5,3.9);
  \node[progress,minimum width=11.8cm] at (7.9,4.65)
    {Validate the terminal QC and enter $(g,t+1)$ with the same generation $g$.\\
     Close term $t$'s missing-proposal test; keep its pending commit timers running.};
  \draw[msg] (13.5,5.45)--node[lab,above] {$\mathrm{PROPOSE}(Y,N_{g,t})$} (7.8,5.9);
  \node[progress,minimum width=11.8cm] at (7.9,6.65)
    {Validate $Y$ and its history; vote immediately in the new term.\\
     Start $Y$'s own $2\Delta$ timer; certification and further proposals continue.};

  \node[phase] at (.1,7.6) {Old and new commit timers overlap};
  \node[state,minimum width=11.8cm] at (7.9,8.35)
    {$T$ commits when its original timer expires, even if the active term is later.\\
     A valid current-generation quit proof first cancels all pending timers in $g$.};
  \node[state,draw=rrAmber,fill=rrAmber!3,minimum width=11.8cm] at (7.9,9.8)
    {Catch-up while $g$ is active: relay a valid current-generation DATA QC's full history.\\
     Start a $2\Delta$ timer if the block has neither a timer nor a local commit.\\
     ANCHOR QCs start no commit timer.};
  \node[anchor=west,align=left,text width=15.4cm,font=\fontsize{9}{10.5}\selectfont] at (.1,11.2)
    {Illustration of the normal path. The diagram does not assert a fixed proposal/QC latency.\\
     No L1 update is needed for this handoff: the generation's leader schedule is already installed.};
\end{tikzpicture}

\caption{Immediate terminal-QC handoff. Certification admits successor DATA votes while existing timers continue. A replica catching up through a DATA QC starts its own timer without re-voting.}
\label{fig:rotation}
\end{figure}

\begin{figure}[htbp]\centering
% Native TikZ fragment. Representative messages; time is not drawn to scale.
\begin{tikzpicture}[
  x=1cm,y=-1cm,font=\fontsize{9.5}{11}\selectfont,
  lane/.style={draw=black!25,densely dashed},
  msg/.style={-{Stealth[length=1.8mm]},semithick,draw=rrBlue},
  lab/.style={fill=white,inner sep=2pt,align=center},
  state/.style={draw=rrBlue!65,fill=white,rounded corners=2pt,
    align=center,inner sep=4pt},
  wait/.style={draw=rrAmber,dashed,fill=rrAmber!4,rounded corners=2pt,
    align=center,inner sep=4pt},
  phase/.style={anchor=west,font=\fontsize{9.5}{11}\selectfont\bfseries}
]
  \path[use as bounding box] (0,-.5) rectangle (15.8,18.0);
  \node[align=center] at (1.5,0) {Finalized L1\\authority record};
  \node[align=center] at (5.7,0) {Correct replica $P$\\active at $(g,t_P)$};
  \node[align=center] at (9.9,0) {Correct replica $Q$\\active at $(g,t_Q)$};
  \node[align=center] at (14.1,0) {New leader\\$L'=L(g+1,0)$};
  \foreach \x in {1.5,5.7,9.9,14.1} \draw[lane] (\x,.6)--(\x,17.15);

  \node[phase] at (.1,.95) {(i) Freeze the current generation once};
  \node[state,draw=rrRed,minimum width=8.0cm] at (8.0,1.7)
    {At $a_P$, $P$ receives valid $\CB(g,t^*)$.\\
     Relay proof; freeze $g$; cancel all its pending timers.};
  \draw[msg,draw=rrRed] (5.7,2.45)--node[lab] {$\CB(g,t^*)$} (9.9,2.9);
  \draw[msg,draw=rrRed] (5.7,2.45)--(14.1,2.9);
  \node[state,draw=rrRed,minimum width=12.6cm] at (7.9,3.65)
    {All correct replicas freeze the same generation, though their terms may differ.\\
     Preserve completed commits; process old messages; produce no new DATA votes.};

  \node[phase] at (.1,4.55) {(ii) Wait for both the drain interval and the next authority record};
  \draw[msg,draw=rrAmber] (1.5,5.2)--node[lab,pos=.54]
    {finalized $\mathrm{AUTH}_{g+1}$; recipients relay it} (14.1,5.65);
  \node[wait,minimum width=12.6cm] at (7.9,6.5)
    {Replica $r$ has the authority record at $d_r$ and enters at
     $e_r=\max\{a_r+\Delta,d_r\}$.\\
     Both readiness spreads are at most $\Delta$, so entry skew is at most $\Delta$.};
  \node[state,minimum width=12.6cm] at (7.9,7.6)
    {At entry, fix highest-QC lock $Q_r$ and send its complete history in STATUS.\\
     Every replica targets $(g+1,0)$; $\mathrm{AUTH}_{g+1}$ uniquely determines $L'$.};
  \draw[msg] (5.7,8.35)--node[lab,pos=.48]
    {$\mathrm{STATUS}(g+1,0,Q_P)$} (14.1,8.75);
  \draw[msg] (9.9,8.95)--node[lab,pos=.52]
    {$\mathrm{STATUS}(g+1,0,Q_Q)$} (14.1,9.4);
  \node[wait,minimum width=8.7cm] at (10.0,10.1)
    {$L'$ waits $2\Delta$ from its own entry $e_{L'}$.\\
     This covers entry skew and correct STATUS delivery.};

  \node[phase] at (.1,10.95) {(iii)--(iv) SH NEWVIEW and anchor re-certification};
  \draw[msg] (14.1,11.6)--node[lab,pos=.53]
    {$\mathrm{NEWVIEW}(g+1,0,Q^*)$\\highest valid received QC} (5.7,12.05);
  \node[state,minimum width=12.6cm] at (7.9,12.85)
    {Validate $Q^*$ and its history; require $\operatorname{rank}(Q^*)\ge\operatorname{rank}(Q_r)$.\\
     Check NEWVIEW equivocation; relay NEWVIEW and vote for its unchanged anchor $A$.};
  \draw[msg] (5.7,13.7)--node[lab] {ANCHOR votes to all} (14.1,14.1);
  \node[state,draw=rrAmber,minimum width=12.6cm] at (7.9,14.85)
    {Form $\QC_{g+1,0}(A)$ from $f+1$ ANCHOR votes.\\
     Anchor re-certification starts no commit timer and does not itself commit $A$.};
  \draw[msg] (14.1,15.65)--node[lab,pos=.53]
    {$\mathrm{PROPOSE}(Y,\QC_{g+1,0}(A))$} (5.7,16.1);
  \node[state,draw=rrTeal,minimum width=12.6cm] at (7.9,16.8)
    {The fresh child $Y$ starts the DATA pipeline and its $2\Delta$ commit timers.};
  \node[anchor=west,align=left,text width=15.4cm,font=\fontsize{9}{10.5}\selectfont] at (.1,17.65)
    {Illustration with a correct new leader. L1 availability may add an unbounded but finite wait.\\
     There is no leader-selection BA. Old-generation proofs do not restart this recovery.};
\end{tikzpicture}

\caption{Generation-wide exit and SH recovery with authority readiness. The local entry time includes both the exit wait and availability of the same finalized next-generation authority. ANCHOR re-certification has no timer.}
\label{fig:recovery}
\end{figure}

\begin{figure}[htbp]\centering
% Native TikZ fragment. Only the displayed worst-case upper bounds are metric.
\begin{tikzpicture}[
  x=1cm,y=-1cm,font=\fontsize{9.5}{11}\selectfont,
  arrow/.style={-{Stealth[length=1.8mm]},semithick},
  state/.style={draw=rrBlue!65,fill=white,rounded corners=2pt,
    align=center,inner sep=5pt},
  lab/.style={fill=white,inner sep=2pt,align=center},
  phase/.style={anchor=west,font=\fontsize{9.5}{11}\selectfont\bfseries}
]
  \path[use as bounding box] (0,-.4) rectangle (15.8,9.4);
  \node[state,minimum width=14.0cm] at (7.9,.45)
    {Assume: a correct initializer, valid first DATA block $B$, and no leader equivocation in this generation.\\
     Let $s$ be the time the initializer first proposes and votes for $B$.};

  \node[phase] at (.1,1.5) {A whole generation need not remain at one leader for $2\Delta$};
  \draw[arrow,draw=black!65] (1.0,2.4)--(15.0,2.4);
  \foreach \x in {1.2,3.3,7.5,13.8} \draw[black!65] (\x,2.28)--(\x,2.52);
  \node[lab,anchor=north] at (1.2,2.65) {$s$};
  \node[lab,anchor=north] at (3.3,2.65) {$s+\Delta$};
  \node[lab,anchor=north] at (7.5,2.65) {$s+3\Delta$};
  \node[lab,anchor=north] at (13.8,2.65) {$s+6\Delta$};
  \draw[very thick,draw=rrTeal] (1.2,2.05)--(7.5,2.05);
  \node[lab,text=rrTeal] at (4.35,1.8) {protected commit window};

  \node[state,text width=3.7cm] (first) at (2.2,4.2)
    {By $s+\Delta$:\\every correct replica has voted for $B$ or seen its DATA QC.};
  \node[state,text width=3.7cm,draw=rrTeal,fill=rrTeal!4] (commit) at (7.5,4.2)
    {A $2\Delta$ timer is active\\at every correct replica.\\All commit $B$ by $s+3\Delta$.};
  \node[state,text width=3.7cm,draw=rrAmber,fill=rrAmber!4] (blame) at (13.1,4.2)
    {A later term's first honest missing-proposal BLAME\\is no earlier than $s+6\Delta$.};
  \draw[arrow,draw=rrTeal] (first.east)--(commit.west);
  \draw[arrow,draw=rrAmber] (commit.east)--(blame.west);

  \node[state,minimum width=14.0cm] at (7.9,6.0)
    {Terminal QCs may move replicas through several ordinary terms during this window.\\
     Those transitions preserve the timers for $B$; a withholding successor cannot yet collect a quit proof.};
  \node[state,draw=rrAmber,fill=rrAmber!3,minimum width=14.0cm] at (7.9,7.45)
    {The initial term does not attract an honest timeout when its correct leader advances on schedule.\\
     With no equivocation proof, every complete BLAME certificate needs at least one honest complaint.};
  \node[anchor=west,align=left,text width=15.4cm,font=\fontsize{9}{10.5}\selectfont] at (.1,8.95)
    {Local progress lemma: a correct initializer in a generation without leader equivocation.\\
     Each completed term contains $K\ge1$ fresh DATA batches; no lower message-delay bound is used here.};
\end{tikzpicture}

\caption{A correct initializer after generation equivocation stops. All first-block vote or QC timers start by $s+\Delta$ and finish by $s+3\Delta$. A later-term honest BLAME cannot be signed before $s+6\Delta$.}
\label{fig:window}
\end{figure}

\begin{figure}[htbp]\centering
% Native TikZ fragment. Proposer exclusion does not change the voting committee.
\begin{tikzpicture}[
  x=1cm,y=-1cm,font=\fontsize{9.5}{11}\selectfont,
  arrow/.style={-{Stealth[length=1.8mm]},semithick},
  state/.style={draw=rrBlue!65,fill=rrBlue!3,rounded corners=3pt,
    align=center,inner sep=5pt,text width=4.2cm,minimum height=1.35cm},
  lab/.style={fill=white,inner sep=3pt,align=center},
  wide/.style={draw=black!40,fill=black!2,rounded corners=3pt,
    align=center,inner sep=5pt,minimum width=14.0cm}
]
  \path[use as bounding box] (0,-.5) rectangle (15.8,11.3);
  \node[wide] at (7.9,.25)
    {Fixed voting committee: $n=2f+1$, at most $f$ static Byzantine identities.\\
     L1 updates proposer eligibility and future schedules; quorum membership stays fixed.};

  \node[state,draw=rrRed,fill=rrRed!4] (proof) at (2.55,2.4)
    {Expose sequencer $F$\\A correct replica receives and relays a valid equivocation proof.};
  \node[state,draw=rrAmber,fill=rrAmber!4] (lone) at (7.9,2.4)
    {L1 inclusion and finality\\Evidence eventually produces a finalized exclusion decision.};
  \node[state,draw=rrTeal,fill=rrTeal!4] (future) at (13.25,2.4)
    {Future authority records\\There is a finite cutoff $G_F$:\\$F$ is ineligible for $g\ge G_F$.};
  \draw[arrow,draw=rrRed] (proof.east)--(lone.west);
  \draw[arrow,draw=rrTeal] (lone.east)--(future.west);

  \node[wide] (fixedauth) at (7.9,4.4)
    {Previously installed $\mathrm{AUTH}_g$ records remain immutable; exclusion is prospective.\\
     No immediate re-entry or new Byzantine identity replaces the excluded proposer.};
  \draw[arrow,draw=rrTeal] (future.south)--(13.25,3.6)--(13.25,3.8);

  \node[state,text width=5.7cm] (finite) at (3.75,6.7)
    {Finite disruption argument\\After its first exposure, each identity has a finite exclusion cutoff.\\
     Only finitely many generation indices precede that cutoff.};
  \node[state,text width=5.7cm,draw=rrTeal,fill=rrTeal!4] (quiet) at (12.05,6.7)
    {Eventual absence of equivocation freezes\\After finitely many such recoveries, no further current-generation equivocation proof can cancel timers.};
  \draw[arrow,draw=rrBlue] (finite.east)--node[lab] {finite faulty\\identity set} (quiet.west);
  \draw[arrow,draw=rrBlue] (3.75,4.9)--(finite.north);
  \node[wide,draw=rrTeal,fill=rrTeal!3] at (7.9,9.05)
    {Combine with schedule fairness and the honest-initializer progress lemma.\\
     Withholding can still cause timeout recovery; it does not invalidate the protected commit window.};
  \draw[arrow,draw=rrTeal] (quiet.south)--(12.05,8.45);
  \node[anchor=west,align=left,text width=15.4cm,font=\fontsize{9}{10.5}\selectfont] at (.1,10.65)
    {Logical dependency diagram. L1 exclusion and activation must progress independently of L2 commitment.\\
     This is an explicit accountability assumption, not a claim that stake alone prevents Byzantine behavior.};
\end{tikzpicture}

\caption{Prospective PoS exclusion. Finalized authority records remain immutable, while each exposed identity is excluded from all sufficiently late generation schedules. Voting membership stays fixed in this reference model.}
\label{fig:exclusion}
\end{figure}

\FloatBarrier

\section{Safety}
\label{sec:safety}
\begin{lemma}[Common recovery target and bounded entry skew]
\label{lem:entry}
Initial-view entry times within each generation differ by at most
$\Delta$. If the first correct generation-$g$ freeze is at $a$, correct
freeze times satisfy $a_r\in[a,a+\Delta]$. Every successor is $(g+1,0)$
with the same finalized authority, and its entry times have spread at most
$\Delta$, even if the authority becomes available only after freezing.
\end{lemma}
\begin{proof}
The initial skew is assumed. Immediate proof relay gives the freeze
bound for replicas already in $g$. A replica still finishing prior recovery
enters $g$ within $\Delta$ of its first correct entry, hence by
$a+\Delta$, and processes the buffered proof on entry. No correct replica
can have completed generation-$g$ recovery before its first correct freeze.
The inductive freeze bound therefore holds.

For any correct $r,s$, both
$|a_r-a_s|\le\Delta$ and $|d_r(g+1)-d_s(g+1)|\le\Delta$ hold.
Taking coordinatewise maxima preserves that bound, so
$e_r=\max\{a_r+\Delta,d_r(g+1)\}$ has spread at most $\Delta$.
All witnesses target the same generation, and L1 supplies one authority
for it. Processing a buffered next-generation proof may immediately
freeze that generation but does not erase its authenticated entry event.
\end{proof}

\begin{lemma}[Canonical history or equivocation]
\label{lem:fork}
Two admitted histories in a generation have one opening and are
prefix-comparable, or their signed material yields an admissible
generation-equivocation proof.
\end{lemma}
\begin{proof}
Distinct openings supply two NEWVIEW signatures, even if the opening
blocks are comparable. Generation zero has a fixed genesis opening.
For the same opening, each fresh height belongs to one term and its
L1-scheduled leader. The first divergent height therefore supplies two
different signed DATA proposals at one $(g,t,k)$. Both histories contain
the terminal-QC admission for that position. Prior-generation evidence
is checked for validity without converting an old fork into a new-generation
quit proof.
\end{proof}

\begin{lemma}[Protection of a successful DATA timer]
\label{lem:timer}
If a correct vote-start or catch-up timer for $B$ starts at $s$ and
commits at $s+2\Delta$, every correct replica knows its DATA QC by expiry.
No admissible QC in that generation can conflict with $B$, including
one assembled later from old votes. Nor can an ANCHOR QC under a
different opening of that generation exist.
\end{lemma}
\begin{proof}
A correct freeze by $s+\Delta$ would relay its proof to the committer
by $s+2\Delta$ and cancel the timer. Therefore no such freeze occurs.
The complete history relayed at timer start reaches all correct replicas
by $s+\Delta$; \cref{lem:entry} ensures that any replica still finishing
prior recovery has entered the generation by then.

For a QC-start timer, the relayed QC itself is known everywhere by
$s+\Delta$. For a vote-start timer, a replica already beyond $B$ has its
QC in the admitted ancestry and relays it; one such relay supplies it
everywhere by expiry. Otherwise all correct replicas vote by $s+\Delta$,
and their at least $f+1$ votes supply the QC by $s+2\Delta$.
An incompatible retained history would instead give an earlier freeze
by \cref{lem:fork}, contradicting successful commitment.

Any conflicting QC has a correct signer. If that signer processes
$B$'s complete history first, comparison prevents its conflicting vote.
Otherwise the conflicting vote and its forwarded signed history precede
that processing, hence occur no later than $s+\Delta$. They reach the
committer by expiry and expose the fork. This contradiction concerns the
time of a correct vote, regardless of when the QC is later assembled.
An ANCHOR vote for a different opening yields the same contradiction
using the two signed NEWVIEWs. Evidence-first processing covers equality.
\end{proof}

\begin{lemma}[Entry locks protect completed timers]
\label{lem:locks}
For a completed DATA timer's block $B$, every correct entry lock in a
later generation ranks at least as high as the protecting QC.
\end{lemma}
\begin{proof}
Let its start be $s$ and the first correct freeze be $a$.
\Cref{lem:timer}'s cancellation argument gives $a>s+\Delta$.
Every successor snapshot is taken at
$e_r\ge a_r+\Delta\ge a+\Delta>s+2\Delta$, after that QC is known
everywhere. Highest-QC selection gives the rank bound, and retained locks
never decrease across later recoveries. Additional L1 waiting only delays
the snapshot. If the generation never freezes, the claim is vacuous.
\end{proof}

\begin{lemma}[Unique extensibility]
\label{lem:extend}
After a correct timer commits $B$, every admitted QC of rank at least
$\rank(QC(B))$ certifies a block extending $B$.
\end{lemma}
\begin{proof}
In the same generation, \cref{lem:timer} excludes alternative certified
openings and incompatible QCs. The one opening's canonical height ranges
increase with the term index; a sufficiently ranked QC cannot certify
a strict ancestor of $B$. It therefore extends $B$.

Suppose a later-generation counterexample exists and choose one of
least rank whose block does not extend $B$. An ANCHOR QC has a correct
voter whose older justification ranks at least as high as its entry lock,
and hence as $QC(B)$ by \cref{lem:locks}. That justification certifies the
same nonextending anchor and has strictly smaller rank, a contradiction.
A DATA QC has a lower-ranked parent QC in that same later generation,
including at a normal boundary or the first fresh child. The parent
cannot extend $B$ when its child does not. It is thus a strictly smaller
counterexample. Lexicographic integer ranks are well-ordered; equal-rank
digest ties do not affect either strict descent.
\end{proof}

\begin{theorem}[SMR safety and validity]
\label{thm:safety}
Under authenticated synchrony and unique finalized authority, all correct
committed logs are prefix-comparable and never withdrawn. Every committed
block is application-valid. These conclusions require neither exclusion,
schedule fairness, nor a positive message-delay lower bound.
\end{theorem}
\begin{proof}
Order the protecting QCs of two direct timer commitments by rank.
\Cref{lem:extend} makes the higher-ranked block extend the lower-ranked
one. This applies equally to vote-start and QC-start timers. Commitments
of ancestors inherit comparability. There is no additional QC-only
instantaneous commit rule. Validated histories supply application
validity, and the protocol cancels unfinished timers without revoking
completed commitments. The authority and signing domain keep entry
identities distinct across generations.
\end{proof}

\section{Honest windows and accountable liveness}
\label{sec:liveness}
\begin{lemma}[Correct-leader progress envelope]
\label{lem:progress}
While a correct term remains unfrozen, by its local deadline
$(2p+4)\Delta$ each correct replica has $p$ admitted fresh positions
supported by own DATA votes or valid DATA QCs, for $p=1,\ldots,K$.
Consequently no correct replica generates a missing-progress BLAME
against such a leader.
\end{lemma}
\begin{proof}
For recovered term zero, let $e$ be earliest entry. The leader enters
by $e+\Delta$ and completes its wait by $e+3\Delta$, after all correct
STATUS messages have arrived. It chooses a QC at least as high as every
fixed entry lock. Its one NEWVIEW reaches all by $e+4\Delta$,
ANCHOR votes give the leader a QC by $e+5\Delta$, and the first fresh
proposal is processed by $e+6\Delta$. Each later proposal/QC iteration
takes at most $2\Delta$. A replica either votes or has QC-backed
catch-up evidence, so the $p$th position is counted by
$e+(2p+4)\Delta$, no later than its own deadline.

For a normal term, first correct receipt of its admitting terminal QC
relays the complete history within $\Delta$. Unless the generation freezes,
all correct entries have that skew. A correct successor's first proposal
is processed within $2\Delta$ of earliest entry; subsequent iterations
again cost at most $2\Delta$. These are within the same generous deadlines.
Completed-term catch-up counts the certified positions and disarms its
old listener. Boundary messages precede timeout checks. Generation zero
starts from genesis and is faster than recovered term zero.
\end{proof}

\Needspace{8\baselineskip}
\begin{lemma}[An uninterrupted DATA window commits]
\label{lem:window}
Suppose a correct replica starts a DATA timer for $B$ and relays its
complete proposal or DATA QC at $s$. If the generation remains unfrozen
at all correct replicas through $s+3\Delta$, every correct replica
commits $B$ by that time. The originating timer commits at $s+2\Delta$.
\end{lemma}
\begin{proof}
Every correct replica processes the complete history by $s+\Delta$.
If it has not passed $B$, it votes and starts its timer; if it has passed
$B$, the admitted ancestry contains its DATA QC, whose processing starts
the catch-up timer when none exists. A prior timer started even earlier,
and an already committed block requires no new action.
Thus all relevant starts are no later than $s+\Delta$ and none is
canceled; every replica commits by $s+3\Delta$.
\end{proof}
The local $2\Delta$ interval begins at a DATA vote or QC-timer start,
not merely at view entry. An honest proposal sent at $s$ can take another
$\Delta$ to reach followers. Also, absence of newly elected Byzantine
leaders does not by itself exclude delayed proof of an older
current-generation equivocation. The window premise excludes such freezes.

\begin{lemma}[Only finitely many generations expose equivocation]
\label{lem:finite}
Under eventual permanent exclusion and the finite fixed Byzantine set,
only finitely many generation indices can be named by admissible
leader-equivocation proofs revealed to correct replicas.
\end{lemma}
\begin{proof}
A correct leader signs one opening and one DATA block per slot, so the
identified signer is Byzantine. For every identity that is ever exposed,
its first revealed proof invokes exclusion and supplies a finite cutoff
$G_p$. It cannot be the admitted leader of any slot in a generation
$g\ge G_p$, so any admissible equivocation attributed to it has $g<G_p$.
There are finitely many such identities and finitely many preceding
generation indices. Old proofs can be revealed later but cannot enlarge
this finite set of generations or restart a completed recovery.
\end{proof}
The number of interruptions need not be at most $f$: one identity can
cause several before its L1 exclusion becomes effective. Finiteness,
not an unstated instantaneous-slashing assumption, is what the proof uses.

\begin{lemma}[Correct initialization after equivocation stops]
\label{lem:initializer}
Consider a generation with correct $L(g,0)$ in which no admissible
equivocation proof is revealed. Let $s$ be its first fresh DATA proposal
and local vote time. Every correct replica commits that block by
$s+3\Delta$, even if ordinary QC rotation reaches a Byzantine successor.
\end{lemma}
\begin{proof}
If no freeze occurs through $s+3\Delta$, apply \cref{lem:window}.
Otherwise take the first correct freeze $a\le s+3\Delta$.
It must use a BLAME certificate because equivocation is excluded.
Term zero has no correct BLAME signature by \cref{lem:progress};
the $f$ Byzantine signatures alone are insufficient.

For a later term $t>0$, choose a correct BLAME signer. It entered that
term at a time $x\ge s$: admission contains the first term's terminal
QC and signed chain, which includes the correct initializer's proposal
created at $s$. Its earliest possible complaint is at $x+6\Delta$.
Consequently any such complete BLAME certificate, including a withheld
one, can first freeze a correct replica only at
$a\ge x+6\Delta\ge s+6\Delta$, contradicting $a\le s+3\Delta$.
Thus the DATA window survives. Every correct replica either votes for
the first block or learns its DATA QC through catch-up by $s+\Delta$;
the two timer rules cover both cases.
\end{proof}

\begin{theorem}[Eventual committed-height progress]
\label{thm:height}
Under \cref{ass:network,ass:l1,ass:exclusion,ass:fairness}, every correct
replica eventually commits arbitrarily large heights with immediate
terminal-QC rotation.
\end{theorem}
\begin{proof}
If infinitely many generations are entered, \cref{lem:finite} leaves
a suffix of generations without revealed equivocation. Initializer
fairness supplies infinitely many correct initial leaders in that suffix.
Each commits its first fresh block at all correct replicas by
\cref{lem:initializer}. Its opening extends every older committed block
by locks and unique extensibility, so its fresh child has strictly greater
height. L1 authority is eventually available for each recovery; the
protocol cannot remain permanently waiting for one next record.

If only finitely many generations are entered, the last generation
cannot freeze, because its wait and next authority readiness both finish.
If only finitely many terms were admitted there, all correct replicas
would eventually catch up to a last common term. If any correct replica
counts all $K$ positions, it either already knows the terminal QC or
DATA-voted the terminal proposal and relays it. With no freeze, every
correct recipient votes or supplies that QC through certified catch-up;
the terminal QC then admits another term. Otherwise every correct
replica fails a scheduled progress check, at latest the $p=K$ check,
and their matching BLAME signatures force recovery. Both contradict
a permanently last term in an unfrozen generation.

Thus normal terms and certified heights increase without bound in that
last generation. Complete proposal/QC relay gives each correct replica
the admitted DATA QCs. Vote or catch-up timers then commit them within
$2\Delta$ of their first starts, since the generation never freezes.
This establishes progress at every correct replica in both cases.
\end{proof}

\begin{theorem}[Request liveness]
\label{thm:requests}
Under the same assumptions, every persistently admissible correct-client
request is eventually committed by all correct replicas.
\end{theorem}
\begin{proof}
With infinitely many generations, infinitely many correct initializers
after the last equivocation generation commit their first fresh blocks
at all correct replicas. With finitely many generations, term fairness
gives infinitely many correct leaders in the final unfrozen generation,
whose blocks eventually commit.

Fix a request after dissemination to all correct replicas. Only finitely
many requests precede it in the common admission order. Each of the
above successful correct proposals includes the oldest fitting pending
requests, or commits them as ancestors. Speculative inclusion does not
remove a request from the queue. Induction over the finite predecessors
therefore commits the request, and safety preserves its position.
\end{proof}

\section{Delay lower bounds, quota, and optional pacing}
\label{sec:pacing}
The $K=1$ perpetual-equivocation execution from the fixed-schedule model
does not apply under \cref{ass:exclusion}: the same identity cannot retain
future proposer eligibility forever after public evidence. Increasing
$K$ and bounding actual delays are a separate way to seek a favorable
commit window; they are not needed for the eventual-exclusion theorem.

A positive physical lower bound $\underline\delta$ alone does not justify
a fresh $2\underline\delta$ certification delay per block. An $f+1$ QC can
contain the correct leader's own vote and all $f$ Byzantine votes.
For a predictable sequence of blocks, Byzantine votes can be sent in
advance and buffered until the leader produces the signed proposals.
The resulting QC can be assembled locally without a fresh round trip
after that proposal. Wire delay and causal certification delay must be
modeled separately.

\paragraph{An explicit optional profile.}
Fix a public $\rho>0$. A correct leader releases its first fresh proposal
as soon as its parent justification is ready. For each subsequent fresh
position, let $q_{k-1}$ be its local parent-QC availability time and
$s_{k-1}$ the previous proposal's first release time. It signs and first
exposes the new proposal at
\[
 s_k=\max\{q_{k-1},s_{k-1}+\rho\}.
\]
It cannot expose the signed proposal earlier in a QC, history bundle, or
other message. It immediately self-processes the released proposal.
All replicas use the profile's progress deadlines
\[
 \operatorname{entry}(v)+6\Delta+
        (p-1)\max\{2\Delta,\rho\},\qquad p=1,\ldots,K.
\]
These policy and deadline changes are part of this optional profile,
not changes to the default immediate-QC protocol.

\begin{proposition}[Optional paced honest-initializer window]
\label{prop:pacing}
In the optional profile just defined, let a correct initial leader
release and DATA-vote its first fresh proposal at $s$. If
$(K-1)\rho>2\Delta$, its first block commits locally before any
ordinary successor is admissible. If $(K-1)\rho>3\Delta$, every correct
replica commits that first block before ordinary handoff.
\end{proposition}
\begin{proof}
The terminal proposal cannot be signed or exposed before $s+(K-1)\rho$.
Its QC and any successor admission require that signed proposal, even
when Byzantine votes are pre-positioned. Before that time, only the
correct initial term is admissible and supplies no equivocation.
The first proposal is processed by earliest entry plus $6\Delta$;
each subsequent proposal/QC step takes at most
$\max\{2\Delta,\rho\}$. Thus every correct progress count meets the
profile's stated deadlines, and no complete BLAME can form.
Old-generation evidence cannot quit this generation. The first
$2\Delta$ timer therefore survives under the first inequality; under
the second inequality \cref{lem:window} applies to every replica.
\end{proof}
For $0<\rho\le2\Delta$, this profile preserves SH's
$2\Delta$ worst-case iteration envelope and original progress deadlines.
This optional profile enforces time spacing explicitly. The default
v3 protocol remains immediate-QC and unpaced; it does not infer these
inequalities from $K\ge1$, $K=n$, or an empirical average delay.

\section{Complexity and latency}
\label{sec:complexity}
Let $\sigma$ be a signature's size, $b$ the fresh payload size, and $A$
the additional history/authority bytes actually transmitted. With semantic
deduplication and full forwarding:
\begin{center}\small
\begin{tabularx}{\linewidth}{@{}>{\raggedright\arraybackslash}p{.21\linewidth}>{\raggedright\arraybackslash}p{.18\linewidth}Y@{}}
\toprule
Operation & Correct messages & Conservative byte cost\\
\midrule
Fresh DATA block & $O(n^2)$ &
$O(n^2(b+A+n\sigma))$ including explicit QC/history forwarding.\\
Normal QC handoff & No extra voting phase &
Admission travels with terminal QC and the successor proposal. No DONE
messages or commitment wait.\\
Catch-up timer & No new vote phase &
Local timer; the first relevant DATA QC/history relay is charged to
that block's dissemination.\\
Progress complaints & $O(n^2)$ per term &
At most one correct BLAME share per replica in an admitted term;
$O(n^2\sigma)$ share bytes plus view/context metadata.\\
Post-freeze recovery & $O(n^2)$ per recovery &
$O(n^2(A+n\sigma))$: quit evidence, STATUS to one leader, NEWVIEW
forwarding, and ANCHOR votes.\\
L1 accountability & Separate control cost &
Evidence submission, finality, epoch activation, and authority proofs;
their cost and delay are not hidden in the local-message count.\\
\bottomrule
\end{tabularx}
\end{center}
Late DATA/QC backlog processed while frozen is charged to its originating
blocks. Complaints that never assemble a quit proof are charged to their
terms, not hidden in a per-generation constant. For normally completed
$K$-block terms, their message cost amortizes to $O(n^2/K)$ per block.
Bounded admitted traffic, not arbitrary Byzantine flooding, is
counted. Explicit signatures and growing histories prevent treating a
QC or admission bundle as a constant-size object without further
compression assumptions.

\paragraph{Normal performance.}
An ordinary DATA vote retains SH's $2\Delta$ local commit timer.
A QC-only catch-up replica waits $2\Delta$ from its first relevant
DATA-QC processing. Normal leader changes preserve both timer types.
In a homogeneous good-case execution with actual delay $\delta$, a leader
can form a parent QC and propose its child in about $2\delta$; terminal-QC
handoff adds no protocol-mandated $\Delta$ wait. The resulting block
interval is network-driven, with a usual $1/(2\delta)$ pipeline rate
under that illustrative timing model. This is not a lower bound on every
possible quorum-formation time or a bandwidth-independent throughput claim.

\paragraph{Recovery with L1 readiness.}
Let $a$ be earliest freeze and $d=\min_r d_r(g+1)$, and define
$E=\max\{a+\Delta,d\}$. Correct successor entries lie in $[E,E+\Delta]$.
Relative to $e=\min_r e_r$, a correct new leader sends NEWVIEW by
$e+3\Delta$, all replicas receive it by $e+4\Delta$, and the leader
has its ANCHOR QC by $e+5\Delta$. The first fresh proposal is processed
by all by $e+6\Delta$; an uninterrupted window gives commitment by
$e+8\Delta$. When the authority was ready in advance, $e=a+\Delta$
and this reduces to the familiar $a+9\Delta$ bound. Otherwise L1
availability adds waiting. A uniform worst-case end-to-end recovery
bound additionally needs a uniform L1 availability bound.

\paragraph{Eventual versus bounded liveness.}
Permanent exclusion and authority availability are eventual assumptions.
They establish eventual SMR progress, not a numerical upper bound on the
time to the last equivocation or the number of aborted generations.
Such bounds require quantitative limits on L1 processing/activation and
on schedule fairness. Neither a slashing penalty nor a nominal epoch
duration alone supplies those limits.

\section{Changes from SH and from v2}
\label{sec:diff}
\begin{center}\small
\begin{tabularx}{\linewidth}{@{}p{.21\linewidth}Y Y@{}}
\toprule
Aspect & Original SH / v2 & v3\\
\midrule
Leader authority & SH: scalar view schedule. V2: preinstalled $L(g,t)$. &
One finalized authority per generation; future schedules implement
PoS exclusion.\\
Normal handoff & SH: stable leader. V2: terminal DONE before new votes. &
Terminal DATA QC immediately admits the next term and its votes.\\
Finality timer & SH/v2: first DATA vote starts $2\Delta$. &
Same fast path; QC-only catch-up also starts $2\Delta$ without re-voting.\\
Quit scope & SH: current view. V2: whole generation. &
Whole generation, including every catch-up timer.\\
Recovery target & V2: $(g+1,0)$. &
Same; wait for its unique authority if not already available.\\
Anchor recovery & SH/v2: lock, NEWVIEW, ANCHOR re-vote without timer. &
Retained; no BA-selected anchor and no automatic anchor commitment.\\
Liveness premise & V2: terminal commitment gates normal admission. &
Finite permanent exclusion of revealed equivocators plus explicit
schedule/input fairness.\\
\bottomrule
\end{tabularx}
\end{center}
Safety does not depend on the slashing service succeeding. The stronger
eventual liveness conclusion does. Without permanent exclusion or with
an indefinitely replenished adversarial identity set, repeated equivocation
can keep canceling timers. Without eventual authority availability,
recovery can instead wait forever for its next finalized record.
Highest-QC inheritance alone resolves neither failure mode.

\begin{thebibliography}{2}\small
\bibitem{sync}
I. Abraham, D. Malkhi, K. Nayak, L. Ren, and M. Yin.
\emph{Sync HotStuff: Simple and Practical Synchronous State Machine
Replication}. IEEE Symposium on Security and Privacy, 2020, Section II,
Figures 1--2.
\href{https://conferences.computer.org/sp/pdfs/sp/2020/349700a106.pdf}{Published paper}.
\bibitem{schneider}
F. B. Schneider. Implementing fault-tolerant services using the state
machine approach: A tutorial. \emph{ACM Computing Surveys} 22(4),
299--319, 1990. \href{https://doi.org/10.1145/98163.98167}{DOI}.
\end{thebibliography}
\end{document}
